%=========== ΛΥΣΗ - 97ο ΘΕΜΑ Γ ===========

\begin{thema}{Γ}

\begin{center}
\begin{tikzpicture}[scale=0.65]
  \draw[help lines, gray!30] (-2,-4) grid (6,4);
  \draw[-latex] (-2.3,0) -- (6.3,0) node[right] {$x$};
  \draw[-latex] (0,-4.3) -- (0,4.3) node[above] {$y$};
  \foreach \x in {-1,1,2,3,4,5}
    \draw (\x,0.1) -- (\x,-0.1) node[below,font=\footnotesize] {$\x$};
  \foreach \y in {-3,-2,-1,1,2,3}
    \draw (0.1,\y) -- (-0.1,\y) node[left,font=\footnotesize] {$\y$};

  \draw (2,-4) -- (2,4);
  \node[font=\footnotesize] at (2.8,3.5) {$x=2$};

  \draw[very thick, maincolor, smooth, tension=0.5]
    plot coordinates {(-1,-1) (0,0) (1,1.5) (1.5,3) (1.8,4)};
  \draw[very thick, maincolor, smooth, tension=0.5]
    plot coordinates {(2.2,-4) (2.3,-2.5) (3,-1) (4,0) (5,0.5)};

  \filldraw[maincolor] (0,0) circle (2pt);
  \filldraw[maincolor] (4,0) circle (2pt);
  \node[maincolor] at (5,2) {$C_f$};
\end{tikzpicture}
\end{center}

\begin{erwthma}
\item Από τη γραφική παράσταση έχουμε
\[
\lim_{x\to 2^-}f(x)=+\infty
\]
και
\[
\lim_{x\to 2^+}f(x)=-\infty.
\]

Θέτουμε
\[
F(x)=\frac{e^{f(x)}-1}{e^{f(x)}+1}.
\]

Όταν $x\to 2^-$, έχουμε $f(x)\to+\infty$, άρα
\[
e^{f(x)}\to+\infty.
\]
Επομένως,
\[
\lim_{x\to 2^-}\frac{e^{f(x)}-1}{e^{f(x)}+1}=1.
\]

Όταν $x\to 2^+$, έχουμε $f(x)\to-\infty$, άρα
\[
e^{f(x)}\to 0.
\]
Επομένως,
\[
\lim_{x\to 2^+}\frac{e^{f(x)}-1}{e^{f(x)}+1}
=\frac{0-1}{0+1}=-1.
\]

Τα δύο πλευρικά όρια είναι διαφορετικά. Άρα το όριο
\[
\lim_{x\to 2}\frac{e^{f(x)}-1}{e^{f(x)}+1}
\]
δεν υπάρχει.

\item Η $f$ είναι παραγωγίσιμη στο $(-1,2)$, άρα είναι συνεχής στο $[0,1]$ και παραγωγίσιμη στο $(0,1)$.

Από το σχήμα έχουμε
\[
f(0)=0
\quad \text{και}\quad
f(1)=\frac{3}{2}.
\]
Άρα, από το Θεώρημα Μέσης Τιμής στο $[0,1]$, υπάρχει $\xi_1\in(0,1)$ τέτοιο ώστε
\[
f'(\xi_1)=\frac{f(1)-f(0)}{1-0}
=\frac{3}{2}.
\]

Επίσης, η $f$ είναι παραγωγίσιμη στο $(2,5)$, άρα είναι συνεχής στο $[3,4]$ και παραγωγίσιμη στο $(3,4)$.

Από το σχήμα έχουμε
\[
f(3)=-1
\quad \text{και}\quad
f(4)=0.
\]
Άρα, από το Θεώρημα Μέσης Τιμής στο $[3,4]$, υπάρχει $\xi_2\in(3,4)$ τέτοιο ώστε
\[
f'(\xi_2)=\frac{f(4)-f(3)}{4-3}
=1.
\]

Επομένως,
\[
f'(\xi_1)\cdot f'(\xi_2)=\frac{3}{2}\cdot 1=\frac{3}{2}.
\]

\item Θεωρούμε τη συνάρτηση
\[
g(x)=f(x)-\frac{1}{x-2},\quad x\in(2,5].
\]
Η εξίσωση
\[
f(x)=\frac{1}{x-2}
\]
είναι ισοδύναμη με
\[
g(x)=0.
\]

Στο διάστημα $[4,5]$ η $g$ είναι συνεχής. Από το σχήμα έχουμε
\[
f(4)=0
\quad \text{και}\quad
f(5)=\frac{1}{2}.
\]
Άρα
\[
g(4)=f(4)-\frac{1}{4-2}=0-\frac{1}{2}=-\frac{1}{2}<0
\]
και
\[
g(5)=f(5)-\frac{1}{5-2}=\frac{1}{2}-\frac{1}{3}=\frac{1}{6}>0.
\]

Επομένως, από το Θεώρημα \eng{Bolzano}, υπάρχει τουλάχιστον μία ρίζα της $g$ στο $(4,5)$, άρα και στο $(2,5)$.

Για τη μοναδικότητα, παρατηρούμε από το σχήμα ότι η $f$ είναι γνησίως αύξουσα στο $(2,5)$. Επίσης, η συνάρτηση
\[
-\frac{1}{x-2}
\]
είναι γνησίως αύξουσα στο $(2,5)$. Άρα η
\[
g(x)=f(x)-\frac{1}{x-2}
\]
είναι γνησίως αύξουσα στο $(2,5)$.

Συνεπώς, η εξίσωση $g(x)=0$ έχει το πολύ μία ρίζα στο $(2,5)$.

Άρα η εξίσωση
\[
f(x)=\frac{1}{x-2}
\]
έχει ακριβώς μία ρίζα στο $(2,5)$.

\item Στο διάστημα $[0,1]$ η $f$ είναι κυρτή. Επομένως η γραφική της παράσταση βρίσκεται κάτω από τη χορδή που ενώνει τα σημεία
\[
(0,f(0))=(0,0)
\]
και
\[
\left(1,f(1)\right)=\left(1,\frac{3}{2}\right).
\]

Η χορδή αυτή έχει εξίσωση
\[
y=\frac{3}{2}x.
\]
Από το σχήμα, η $C_f$ δεν ταυτίζεται με τη χορδή στο εσωτερικό του διαστήματος $[0,1]$. Άρα, για κάθε $x\in(0,1)$,
\[
f(x)<\frac{3}{2}x.
\]

Ολοκληρώνοντας στο $[0,1]$, παίρνουμε
\[
\int_0^1 f(x)\,\d x
<
\int_0^1 \frac{3}{2}x\,\d x
=\frac{3}{2}\cdot \frac{1}{2}
=\frac{3}{4}.
\]
\end{erwthma}
\end{thema}

%=========== ΤΕΛΟΣ ΛΥΣΗΣ - 97ο ΘΕΜΑ Γ ===========
